\newpage
\section*{CHƯƠNG 4. THỰC NGHIỆM VÀ ĐÁNH GIÁ  }
\addcontentsline{toc}{section}{\numberline{}CHƯƠNG 4. THỰC NGHIỆM VÀ ĐÁNH GIÁ}


\setcounter{section}{4}
\setcounter{subsection}{0}
\setcounter{figure}{0}
\setcounter{table}{0}

\subsection{Mục tiêu thực nghiệm}

Chương này đặt mục tiêu thiết lập hệ thống thực nghiệm nhằm đánh giá toàn diện hiệu quả, độ ổn định và khả năng tổng quát hóa của mô hình \textbf{MaRENet} trong bối cảnh phát hiện gian lận thẻ tín dụng theo chuỗi thời gian. Thực nghiệm được thiết kế xoay quanh ba trục phân tích chính.

Trước hết, chương trình thực nghiệm tiến hành đánh giá \textit{hiệu quả định lượng} của MaRENet so với các mô hình nền phổ biến trên hai tập dữ liệu \textit{Main} và \textit{External}, được mô tả chi tiết trong mục ~\hyperref[sec:dataset]{\texttt{Bộ dữ liệu}}. Việc so sánh dựa trên các chỉ số nhạy với lớp thiểu số như Precision, Recall và F1-score, qua đó làm rõ khả năng xử lý bài toán mất cân bằng dữ liệu nghiêm trọng.

Tiếp theo, quá trình thực nghiệm kiểm định \textit{khả năng tổng quát hóa} của mô hình khi chuyển miền dữ liệu. Sự khác biệt về thời gian và kênh giao dịch giữa tập \textit{Main} và \textit{External} cho phép kiểm tra mức độ độc lập của kiến trúc đối với phân phối huấn luyện ban đầu, thông qua việc theo dõi sự suy giảm của các chỉ số đánh giá.

Sau cùng, luận văn tiến hành phân tích \textit{đóng góp từng thành phần} trong kiến trúc thông qua các thí nghiệm ablation. Việc loại bỏ có kiểm soát từng thành phần như Feature-wise Mask, Focal Loss và khối RNN cho phép xác định mức độ ảnh hưởng và vai trò thiết kế của từng module trong tổng thể mô hình.

Ba hướng phân tích trên đóng vai trò trụ cột trong việc xác thực hiệu năng và tính ứng dụng của MaRENet trong môi trường tài chính thực tế, nơi yêu cầu cao về tính chính xác và độ phản ứng thời gian thực.

\subsection{Thiết lập thực nghiệm}

\subsubsection{Môi trường tính toán.}

Toàn bộ thực nghiệm được thực hiện bằng ngôn ngữ lập trình \texttt{Python\,3.10}, với khung học sâu \texttt{PyTorch\,2.2.0}. Hệ thống máy trạm sử dụng bộ xử lý \textbf{AMD Ryzen Threadripper PRO 5975WX} (32 nhân, 64 luồng, xung nhịp tối đa 7.0 GHz), bộ nhớ RAM vật lý 256 GB, trong đó có 202 GB khả dụng. Quá trình huấn luyện và suy luận mô hình được tăng tốc bởi hai GPU \texttt{NVIDIA GeForce RTX 4090}, mỗi card có 24 GB VRAM, được quản lý qua trình điều khiển \texttt{NVIDIA-SMI 565.77} và hỗ trợ \texttt{CUDA 12.7}.

Đối với các mô hình học máy truyền thống, thư viện \texttt{scikit-learn\,1.5.0} được sử dụng để triển khai Logistic Regression, Decision Tree và Random Forest. Hai mô hình boosting hiện đại là \texttt{LightGBM\,4.3.0} và \texttt{CatBoost\,1.2.2} cũng được tích hợp nhằm mở rộng phổ so sánh. Toàn bộ quá trình huấn luyện, theo dõi chỉ số, lưu checkpoint và trực quan hóa biểu đồ được tổ chức thông qua nền tảng \texttt{Weights \& Biases} (wandb).

Hệ điều hành sử dụng trong thí nghiệm là \texttt{Ubuntu 22.04 LTS}, môi trường ảo được quản lý bằng \texttt{venv}, đảm bảo tính tái lập và độc lập cho mỗi lần chạy mô hình.

\subsubsection{Chiến lược huấn luyện.}

Trong quá trình huấn luyện, mô hình MaRENet được tối ưu hóa bằng thuật toán \textbf{Adam} với tốc độ học (\textit{learning rate}) mặc định là \(10^{-3}\), một lựa chọn phổ biến giúp mô hình hội tụ nhanh và ổn định trong các bài toán chuỗi thời gian. Kích thước lô (\textit{batch size}) được đặt là 256 nhằm cân bằng giữa tốc độ huấn luyện và tính chính xác của gradient.

Mô hình được huấn luyện tối đa trong 100 vòng lặp (\textit{epochs}). Tuy nhiên, để tránh hiện tượng quá khớp (\textit{overfitting}) và tối ưu hiệu quả thời gian, chiến lược \textbf{early stopping} được áp dụng. Cụ thể, quá trình huấn luyện sẽ tự động kết thúc nếu sau 8 epoch liên tiếp không ghi nhận sự cải thiện nào về chỉ số F1-score trên tập kiểm tra. Điều này đảm bảo rằng mô hình không chỉ học tốt trên tập huấn luyện mà còn duy trì được hiệu suất cao trên dữ liệu chưa từng thấy.


\subsubsection{Bộ dữ liệu.}
\label{sec:dataset}

Các thí nghiệm trong nghiên cứu được triển khai trên bộ dữ liệu \textit{Credit Card Transactions Fraud Detection}~\cite{KaggleFraudData}, bao gồm khoảng 1{.}85 triệu giao dịch thẻ tín dụng thực tế được thu thập từ 999 người dùng khác nhau. Mỗi giao dịch được mô tả bởi 22 thuộc tính, bao phủ nhiều loại thông tin như thời gian giao dịch, đặc trưng hành vi và nhân khẩu học của chủ thẻ. Các trường dữ liệu bao gồm cả dạng định lượng (như số tiền giao dịch) và định danh dạng rời rạc (như mã bưu điện, số thẻ, nghề nghiệp), phản ánh đầy đủ đặc trưng nghiệp vụ trong các hệ thống thanh toán hiện đại.

Mặc dù được xây dựng dưới dạng mô phỏng, bộ dữ liệu này được tạo ra bằng công cụ \textit{Sparkov} kết hợp với thư viện \textit{Faker}, nhằm tái hiện các giao dịch tài chính từ tháng 1/2019 đến tháng 12/2020 với độ chi tiết cao.

Từ các thuộc tính gốc, đặc biệt là \texttt{trans\_date\_trans\_time}, tập đặc trưng được mở rộng với các thông tin thời gian có giá trị cao như ngày trong tuần và giờ trong ngày, giúp mô hình phát hiện các bất thường theo khung giờ, chu kỳ hoặc nhịp độ tiêu dùng. Ngoài ra, các trường như \texttt{amt}, \texttt{cc\_num}, \texttt{category}, \texttt{job}, \texttt{cust\_age}, \texttt{cust\_age\_groups}, \texttt{merchant\_num}, \texttt{zip\_num}, \texttt{state\_num}, v.v. được mã hóa lại để phù hợp với xử lý mô hình hóa, tạo thành một tập dữ liệu giàu ngữ cảnh, phản ánh hành vi người dùng một cách toàn diện.

Về mặt nhãn, chỉ dưới 0{,}6\% các giao dịch được gán là gian lận, tái hiện chính xác đặc trưng mất cân bằng nghiêm trọng trong các hệ thống tài chính thực tế. Điều này đặt ra thách thức lớn cho mô hình học máy, đòi hỏi chiến lược huấn luyện và đánh giá đặc biệt.

Để đảm bảo đánh giá khách quan và tránh rò rỉ thông tin, dữ liệu được phân chia theo người dùng: 80\% người dùng thuộc về tập huấn luyện và 20\% còn lại cho tập kiểm thử. Trong phần huấn luyện, một tập kiểm định chiếm 10\% được tách riêng theo thứ tự thời gian để giám sát quá trình học. Tiền xử lý chỉ thực hiện các bước tối thiểu như loại bỏ bản ghi trùng lặp, hiệu chỉnh dấu thời gian sai định dạng và bù khuyết các giá trị phân loại theo mode nhằm bảo toàn độ chân thực của dữ liệu gốc.

Bên cạnh bộ dữ liệu chính, nghiên cứu còn sử dụng một tập dữ liệu thứ hai có cùng sơ đồ đặc trưng để kiểm tra độ bền mô hình khi phân phối dữ liệu thay đổi. Tập dữ liệu này được thu thập từ nguồn công khai khác trên nền tảng Kaggle~\cite{ExternalFraudData}, với các đặc trưng tương đồng nhưng được tổng hợp từ một tập khách hàng và hành vi khác biệt, nhằm tạo ra một môi trường thử nghiệm đa dạng hơn. Nó đóng vai trò như một tập \textit{External} giúp đánh giá khả năng tổng quát hóa của MaRENet khi áp dụng vào các miền dữ liệu khác biệt với miền huấn luyện. Việc mô hình duy trì hiệu năng ổn định trên tập này sẽ là minh chứng rõ ràng cho độ tin cậy của kiến trúc trong các kịch bản triển khai thực tế.


\subsubsection{Các mô hình so sánh.}

Để đánh giá tính hiệu quả của kiến trúc MaRENet, nghiên cứu triển khai so sánh với hai nhóm baseline phổ biến:

\begin{itemize}
    \item \textbf{Thuật toán học máy truyền thống}: Bao gồm K-Nearest Neighbors (KNN), Logistic Regression (LR), Decision Tree (DT), Random Forest (RF), LightGBM, CatBoost, XGBoost, cũng như hai phương pháp tổ hợp là Stacking Classifier và Voting Classifier. Các thuật toán này đại diện cho phổ rộng của các tiếp cận cổ điển, từ tuyến tính đến cây quyết định và học tăng cường gradient.
    
    \item \textbf{Mô hình học sâu}: Bao gồm GRU và LSTM thuần để đánh giá riêng khả năng học chuỗi. Ngoài ra, hai kiến trúc SOTA được lựa chọn là \textit{STAN ~\cite{cheng2020spatio}} (Spatio-Temporal Attention Network) và \textit{TH-LSTM ~\cite{Xie2022}} (Temporal Hierarchical LSTM), vốn đã chứng minh hiệu quả vượt trội trong phát hiện gian lận theo chuỗi thời gian.
\end{itemize}

Tất cả các mô hình đều được huấn luyện thống nhất trên cùng pipeline tiền xử lý, áp dụng kỹ thuật chuẩn hóa đặc trưng và chia dữ liệu theo người dùng như MaRENet. Việc lựa chọn siêu tham số được thực hiện bằng grid search trên tập validation nhằm đảm bảo đánh giá công bằng và tối ưu hóa từng mô hình theo tiêu chí khách quan.

\vspace{0.5em}
\subsubsection{Tiêu chí đánh giá.}

Do đặc trưng mất cân bằng nghiêm trọng trong phân bố nhãn—với tỷ lệ giao dịch gian lận dưới 0,6\%—việc lựa chọn thước đo đánh giá là yếu tố then chốt trong nghiên cứu. Bốn chỉ số được sử dụng bao gồm:

\begin{itemize}
    \item \textbf{Accuracy} (Độ chính xác tổng thể): Tỷ lệ dự đoán đúng trên toàn bộ tập kiểm thử, tuy nhiên có thể gây hiểu nhầm trong bài toán mất cân bằng.
    \item \textbf{Precision} (Độ chính xác của lớp gian lận): Tỷ lệ giao dịch được dự đoán là gian lận và thực sự là gian lận.
    \item \textbf{Recall} (Tỷ lệ phát hiện gian lận): Tỷ lệ giao dịch gian lận thực sự được mô hình phát hiện ra.
    \item \textbf{F1-score} (Hàm điều hòa Precision và Recall): Là chỉ số chính được sử dụng trong đánh giá vì phản ánh đồng thời năng lực phát hiện đúng (Recall) và hạn chế báo động giả (Precision) trên lớp thiểu số.
\end{itemize}

Các chỉ số này được tính toán dựa trên ma trận nhầm lẫn gồm các thành phần: 

\begin{align}
\text{Accuracy}  &= \frac{TP + TN}{TP + TN + FP + FN} \label{eq:accuracy} \\
\text{Precision} &= \frac{TP}{TP + FP} \label{eq:precision} \\
\text{Recall}    &= \frac{TP}{TP + FN} \label{eq:recall} \\
\text{F1-score}  &= 2 \cdot \frac{\text{Precision} \cdot \text{Recall}}{\text{Precision} + \text{Recall}} \label{eq:f1score}
\end{align}



Trong đó, \textbf{TP} (True Positives) là số giao dịch gian lận được dự đoán đúng; \textbf{FP} (False Positives) là số giao dịch hợp lệ bị dự đoán sai là gian lận; \textbf{FN} (False Negatives) là số giao dịch gian lận bị bỏ sót; và \textbf{TN} (True Negatives) là số giao dịch hợp lệ được dự đoán đúng.

Việc tập trung vào F1-score giúp đánh giá toàn diện năng lực của mô hình trong việc xử lý các tình huống nhạy cảm, nơi sai sót có thể dẫn đến hậu quả nghiêm trọng trong thực tế vận hành hệ thống tài chính.

\subsection{Kết quả đánh giá mô hình}
\begin{table*}[ht]
\centering
\scriptsize
\resizebox{0.95\textwidth}{!}{%
\begin{tabular}{l|cccc|cccc}
\hline
\textbf{Model} & \multicolumn{4}{c|}{\textbf{Main Dataset}} &
                 \multicolumn{4}{c}{\textbf{External Dataset}}\\
               & Accuracy & Precision & Recall & F1~Score &
                 Accuracy & Precision & Recall & F1~Score\\
\hline
KNN ~\cite{cover1967nearest}                    & 0.9965 & 0.7559 & 0.4970 & 0.5997 & 0.9891 & 0.5073 & 0.5356 & 0.5211\\
Logistic Regression ~\cite{cox1958regression}    & 0.9948 & 0.5013 & 0.3623 & 0.4206 & 0.9888 & 0.4927 & 0.3943 & 0.4380\\
Decision Tree ~\cite{quinlan1986induction}          & 0.9965 & 0.7113 & 0.5485 & 0.6194 & 0.9865 & 0.3949 & 0.4094 & 0.4020\\
LightGBM ~\cite{ke2017lightgbm}               & 0.9978 & 0.8069 & 0.7570 & 0.7811 & 0.9902 & 0.5649 & 0.5039 & 0.5326\\
Random Forest ~\cite{breiman2001random}          & 0.9988 & 0.9182 & 0.8383 & 0.8764 & 0.9904 & 0.5816 & 0.4657 & 0.5173\\
CatBoost ~\cite{prokhorenkova2018catboost}               & 0.9987 & 0.9306 & 0.8082 & 0.8651 & 0.9900 & 0.5510 & 0.5132 & 0.5314\\
XGBoost ~\cite{chen2016xgboost}                & 0.9990 & 0.9426 & 0.8612 & 0.9001 & 0.9901 & 0.5586 & 0.5099 & 0.5331\\
Stacking Classifier ~\cite{wolpert1992stacked}    & 0.9989 & 0.9271 & 0.8619 & 0.8933 & 0.9903 & 0.5614 & 0.5416 & 0.5513\\
Voting Classifier ~\cite{zhou2012ensemble}      & 0.9990 & 0.9276 & 0.8691 & 0.8974 & 0.9902 & 0.5614 & 0.5166 & 0.5381\\
GRU ~\cite{cho2014learning}            & 0.9995 & 0.9701 & 0.9286 & 0.9489 & 0.9993 & 0.9626 & 0.9706 & 0.9666\\
LSTM ~\cite{hochreiter1997long}           & 0.9995 & \textbf{0.9782} & 0.9241 & 0.9484 & 0.9994 & 0.9792 & 0.9569 & 0.9679\\
STAN ~\cite{cheng2020spatio}                  & 0.9990 & 0.9184 & 0.8879 & 0.9029 & 0.9989 & 0.9654 & 0.9242 & 0.9443\\
TH\textendash LSTM ~\cite{Xie2022}     & 0.9994 & 0.9522 & 0.9307 & 0.9413 & \underline{0.9994} & \underline{0.9854} & 0.9532 & 0.9690\\
\textbf{MaRENet\textendash GRU}  & \textbf{0.9997} & \underline{0.9765} & \textbf{0.9623} & \textbf{0.9694} &
                               \textbf{0.9995} & 0.9848 & \underline{0.9652} & \underline{0.9749}\\
\textbf{MaRENet\textendash LSTM} & \underline{0.9996} & 0.9730 & \underline{0.9583} & \underline{0.9656} &
                               \textbf{0.9995} & \textbf{0.9865} & \textbf{0.9677} & \textbf{0.9770}\\
\hline
\end{tabular}}
\caption{Hiệu năng so sánh trên hai bộ dữ liệu. Giá trị \textbf{in đậm} là tốt nhất, \underline{gạch chân} là nhì tốt.}
\label{tab:mergeresults}
\end{table*}

Phần thực nghiệm tiến hành so sánh toàn diện giữa các mô hình học máy (ML) và học sâu (DL) trong bối cảnh phát hiện gian lận thẻ tín dụng. Nhóm mô hình học máy bao gồm các thuật toán cổ điển như K-Nearest Neighbors (KNN), Logistic Regression, Decision Tree, cùng với các mô hình tổ hợp hiện đại như LightGBM, Random Forest, CatBoost, XGBoost và hai bộ phân loại Voting và Stacking. Đối với nhóm học sâu, nghiên cứu triển khai các mạng hồi tiếp đơn thuần GRU và LSTM, bên cạnh hai kiến trúc tiên tiến là STAN~\cite{cheng2020spatio} và TH-LSTM~\cite{Xie2022}.

Tất cả mô hình đều được huấn luyện trong cùng một quy trình tiền xử lý, áp dụng chiến lược tinh chỉnh siêu tham số nhất quán trên tập validation. Bốn chỉ số đánh giá gồm Accuracy, Precision, Recall và F1-score được sử dụng, với F1-score đóng vai trò trọng tâm nhằm phản ánh sự cân bằng giữa năng lực phát hiện và độ chính xác trong phân loại giao dịch gian lận.

Bảng~\ref{tab:mergeresults} tổng hợp kết quả thực nghiệm cho thấy rõ ưu thế của hai biến thể đề xuất \emph{MaRENet–GRU} và \emph{MaRENet–LSTM}. Trên \textbf{Main dataset}, MaRENet–GRU đạt F1-score cao nhất (0{.}9694), vượt qua toàn bộ các baseline về mọi mặt, kể cả các mô hình học sâu hiện đại như TH-LSTM hay STAN. Đáng chú ý, MaRENet–GRU thể hiện khả năng phát hiện gian lận mạnh mẽ hơn đáng kể so với XGBoost—mô hình truyền thống đạt kết quả tốt nhất—với F1-score cao hơn tương đối nhiều và Recall được nâng lên rõ rệt, giúp phát hiện thêm nhiều trường hợp gian lận hiếm gặp mà vẫn duy trì Precision ở mức rất cao.

Trên \textbf{External dataset}, MaRENet–LSTM vươn lên dẫn đầu với F1-score đạt 0{,}9770, thể hiện độ ổn định ấn tượng khi chuyển sang miền dữ liệu chưa từng thấy trong huấn luyện. Cả hai biến thể MaRENet chỉ suy giảm F1-score rất nhẹ khi chuyển miền, trong khi phần lớn mô hình truyền thống bị ảnh hưởng rõ rệt, giảm hiệu năng đáng kể.

Một điểm khác biệt thú vị nằm ở chênh lệch Precision và Recall giữa hai biến thể: MaRENet–GRU thiên về bắt thêm các gian lận mới, chấp nhận đánh đổi đôi chút độ chính xác; trong khi MaRENet–LSTM lại kiểm soát cảnh báo giả tốt hơn, nhưng dễ bỏ sót các trường hợp ở ranh giới phân loại. Việc lựa chọn giữa hai kiến trúc có thể linh hoạt tuỳ theo mục tiêu nghiệp vụ: ưu tiên an toàn tối đa, hay nhắm đến khả năng phát hiện toàn diện hơn.

Quan sát kỹ các lỗi phân loại còn sót lại cho thấy phần lớn nằm ở các giao dịch có giá trị nhỏ, diễn ra vào khung giờ quen thuộc với hành vi người dùng – vốn dễ bị mô hình “hiểu lầm” là hợp lệ. Điều này mở ra hướng nghiên cứu kết hợp thêm đặc trưng theo dõi hành vi dài hạn hoặc dữ liệu tín dụng không gian chi tiết hơn.

Tổng kết lại, MaRENet không chỉ chứng minh hiệu quả vượt trội trên cả hai bộ dữ liệu mà còn thể hiện độ bền vững trước thay đổi phân phối – một yêu cầu then chốt đối với hệ thống phát hiện gian lận hoạt động trong môi trường thực tế liên tục biến động.

\subsection{Phân tích ảnh hưởng thành phần (ablation)}
\label{subsec:ablation}

\subsubsection{Ảnh hưởng của Channel-wise Masking và Focal Loss}

Bảng~\ref{tab:ablation-core} làm rõ vai trò thiết yếu của hai thành phần cốt lõi trong kiến trúc MaRENet: khối \emph{channel-wise masking} và hàm mất mát \emph{Focal Loss}. Khi loại bỏ cơ chế masking—vốn giúp mô hình tập trung vào các chiều thông tin quan trọng trong không gian đặc trưng—hiệu năng của mô hình suy giảm rõ rệt. Độ chính xác tổng thể giảm nhẹ, trong khi chỉ số F1 sụt giảm tương đối đáng kể, phản ánh sự suy yếu trong khả năng nhận diện chính xác các mẫu gian lận.

Tương tự, khi thay thế \emph{Focal Loss} bằng hàm mất mát cổ điển như Binary Cross-Entropy, mô hình trở nên kém nhạy hơn với các trường hợp khó phân loại. Điều này dẫn đến sự suy giảm nhất định về Recall và kéo theo hiệu quả tổng thể bị ảnh hưởng ở mức không thể bỏ qua.

Kết quả ablation này khẳng định rằng: cả cơ chế chọn lọc đặc trưng theo chiều và hàm mất mát thích nghi đều đóng vai trò không thể thiếu trong việc duy trì hiệu suất vượt trội của MaRENet trong bài toán phát hiện gian lận.

\begin{table}[H]
\centering
\small
\resizebox{0.9\textwidth}{!}{%
\begin{tabular}{lcccc}
\hline
\textbf{Biến thể} & \textbf{Accuracy} & \textbf{Precision} & \textbf{Recall} & \textbf{F1} \\
\hline
\textbf{MaRENet–GRU (đầy đủ)} & \textbf{0.9997} & 0.9765 & \textbf{0.9623} & \textbf{0.9694} \\
\quad– không \emph{Channel-wise Masking} & 0.9994 & 0.9662 & 0.9186 & 0.9418 \\
\quad– không \emph{Focal Loss} & 0.9996 & \textbf{0.9815} & 0.9337 & 0.9570 \\
\hline
\end{tabular}}
\caption{Ablation về cơ chế \emph{masking} và hàm \emph{Focal Loss}. Cả hai thành phần đều đóng vai trò thiết yếu trong việc duy trì hiệu năng SOTA của MaRENet.}
\label{tab:ablation-core}
\end{table}


\subsubsection{Ảnh hưởng của chiều dài bộ nhớ giao dịch \(d\)}

\begin{table}[H]
\centering
\caption{Ảnh hưởng của chiều dài bộ nhớ \(d\) tới F1-score.}
\label{tab:memory-length}
\begin{tabular}{|c|c|}
\hline
\textbf{Chiều dài \(d\)} & \textbf{F1 Score}  \\
\hline
50   & 0.9568  \\
200  & 0.9600  \\
500  & 0.9627 \\
1854 (heuristic) & \textbf{0.9694} \\
\hline
\end{tabular}
\end{table}

Một yếu tố quan trọng khác là chiều dài bộ nhớ giao dịch \(d\)—tức số lượng giao dịch quá khứ được đưa vào mỗi cửa sổ đầu vào.  
Như trình bày ở Công thức~(4), giá trị \(d = 1854\) được xác định theo thống kê toàn bộ lịch sử giao dịch trung bình mỗi người. Tuy nhiên, trong một số ứng dụng thực tế (ví dụ hệ thống thời gian thực), việc sử dụng toàn bộ lịch sử có thể không khả thi.  

Bảng~\ref{tab:memory-length} cho thấy khi tăng \(d\), hiệu năng mô hình cải thiện nhờ tiếp nhận thêm bối cảnh hành vi. Tuy nhiên, mức cải thiện dần chậm lại sau mốc \(d = 500\), trong khi chi phí tính toán vẫn tăng tuyến tính.  
Giá trị \(d = 1854\) cho kết quả tốt nhất, nhưng các lựa chọn trung gian như \(d = 200\) hay \(d = 500\) vẫn đạt F1 khá cao và có thể được ưu tiên trong môi trường bị giới hạn tài nguyên.

\subsubsection{Ảnh hưởng của phép chuẩn hoá theo ngày/giờ}

Để đánh giá vai trò của thông tin thời gian trong hành vi giao dịch, chúng tôi thực hiện ablation trên module \emph{Day–Hour Aware}—một thành phần học trọng số điều chỉnh cho hai đặc trưng \texttt{day\_of\_week} và \texttt{trans\_hour}. Trong mô hình đầy đủ, hai đặc trưng này không được xử lý bằng các hằng số cố định mà được nhân với các hệ số huấn luyện được, cho phép mô hình điều chỉnh mức độ nhấn mạnh vào các hành vi giao dịch mang tính chu kỳ như chi tiêu cuối tuần hoặc các giao dịch bất thường vào ban đêm.

Khi tắt module này (giữ hai trọng số luôn bằng 1.0), hiệu năng mô hình suy giảm thấy rõ. Cụ thể, Bảng~\ref{tab:ablation-dayhour} cho thấy điểm F1 giảm khoảng 1{,}2 điểm, chứng minh rằng việc khai thác thời gian giao dịch một cách thích ứng giúp nâng cao độ nhạy của hệ thống với các tín hiệu gian lận tinh vi.

\begin{table}[H]
\centering
\small
\begin{tabular}{lc}
\hline
\textbf{Biến thể} & \textbf{F1} \\
\hline
\textbf{MaRENet–GRU (có Day/Hour scaling)} & \textbf{0.9694} \\
– không \emph{Day/Hour scaling} & 0.9576 \\
\hline
\end{tabular}
\caption{Ablation module chuẩn hoá theo ngày/giờ. Việc học trọng số thời gian đem lại lợi ích rõ rệt cho mô hình.}
\label{tab:ablation-dayhour}
\end{table}

\subsection{Phân tích hiệu quả tính toán}
\label{subsec:efficiency}

\paragraph{Đánh giá FLOPs.}
\mbox{}\\
FLOPs (\textit{Floating Point Operations}) phản ánh trực tiếp độ phức tạp tính toán cho một lượt suy luận (forward pass).  
Với MaRENet, một lượt truyền xuôi chỉ tiêu tốn khoảng \textbf{52{,}48 triệu} FLOPs – con số khá nhỏ so với nhiều kiến trúc học sâu hiện hành.  
Nhờ vậy, mô hình \textbf{có tiềm năng đạt độ trễ suy luận thấp} khi triển khai trên GPU hiện đại; tuy nhiên \textbf{việc xác nhận cụ thể độ trễ (latency) vẫn cần được đo lường thực nghiệm} 

\paragraph{Phân tích độ nhẹ (Lightweight Analysis).}
\mbox{}\\
Ngoài chi phí FLOPs khiêm tốn, MaRENet còn sở hữu số tham số nhỏ gọn. Như Bảng~\ref{tab:params} cho thấy, toàn bộ mạng—bao gồm khối \emph{channel-wise masking} và GRU—chỉ có \textbf{52{,}822} tham số huấn luyện, thấp hơn đáng kể so với TH-LSTM và STAN.  
Kích thước gọn nhẹ giúp rút ngắn thời gian huấn luyện, giảm nguy cơ overfitting và \textbf{gợi ý khả năng hoạt động hiệu quả trong các pipeline giám sát gian lận đòi hỏi độ trễ thấp}. Việc benchmark latency thực tế sẽ được xem xét trong nghiên cứu tiếp theo.

\begin{table}[H]
\centering
\small
\begin{tabular}{|l|c|}
\hline
\textbf{Mô hình} & \textbf{Số tham số huấn luyện} \\
\hline
TH-LSTM        & 143{,}745 \\
STAN           & 75{,}593  \\
\textbf{MaRENet (đề xuất)} & \textbf{52{,}822} \\
\hline
\end{tabular}
\caption{So sánh số tham số giữa các mô hình.}
\label{tab:params}
\end{table}

\paragraph{Giải thích đóng góp đặc trưng qua mặt nạ học được.}
\mbox{}\\
MaRENet cung cấp tính minh bạch thông qua cơ chế \emph{channel-wise masking}: mỗi đặc trưng đầu vào được gán trọng số chú ý học được trong quá trình huấn luyện.  
Hình~\ref{fig:mask_values} minh hoạ các giá trị attention sau khi mô hình hội tụ; một nhóm nhỏ đặc trưng—\texttt{category}, \texttt{amt}, \texttt{trans\_hour}, \texttt{is\_night}, nhóm tuổi, mã thương hiệu và giới tính—liên tục giữ trọng số cao, cho thấy vai trò then chốt của tín hiệu ngữ cảnh-thời gian và nhân khẩu học trong phát hiện gian lận.

\begin{figure}[H]
\centering
\includegraphics[width=0.8\linewidth]{Images/figures_mask_values.png}
\caption{Trọng số attention học được trên các đặc trưng đầu vào. Một vài đặc trưng chủ chốt chiếm ưu thế rõ rệt.}
\label{fig:mask_values}
\end{figure}

\subsection{Kết luận chương 4}

Chương 4 đã đánh giá toàn diện hiệu quả và độ ổn định của mô hình \textbf{MaRENet} trong bài toán phát hiện gian lận thẻ tín dụng.  
Trên cả hai tập dữ liệu \textit{Main} và \textit{External}, MaRENet thể hiện độ chính xác vượt trội, vượt hơn đáng kể so với các mô hình nền, bao gồm cả thuật toán truyền thống lẫn các kiến trúc học sâu hiện đại như GRU, LSTM, STAN và TH-LSTM.  
Kết quả này phản ánh khả năng tổng quát hóa mạnh mẽ của MaRENet trước các thay đổi về phân phối dữ liệu, đồng thời khẳng định hiệu quả của thiết kế kiến trúc gọn nhẹ nhưng giàu thông tin.

Bên cạnh đó, việc phân tích đã làm rõ vai trò thiết yếu của từng thành phần trong kiến trúc. Cơ chế \textbf{mặt nạ đặc trưng} cho phép mô hình tập trung chọn lọc vào các chiều đặc trưng quan trọng, giúp giảm nhiễu và nâng cao độ nhạy. Tiếp theo, việc áp dụng hàm mất mát \textit{Focal Loss} đã chứng minh hiệu quả rõ rệt trong việc xử lý mất cân bằng dữ liệu, đặc biệt là cải thiện khả năng phát hiện lớp gian lận hiếm gặp. Hơn nữa, module \textit{Day–Hour Aware} giúp mô hình khai thác tối đa các quy luật hành vi mang tính chu kỳ theo thời gian, từ đó tăng cường khả năng nhận diện các hành vi bất thường. Đồng thời, chiều dài bộ nhớ $d$ cũng cho thấy ảnh hưởng rõ rệt đến hiệu năng, khi việc mở rộng ngữ cảnh giao dịch giúp mô hình nắm bắt tốt hơn các chuỗi hành vi gian lận tiềm ẩn.


Về mặt hiệu suất tính toán, MaRENet duy trì mức độ tính toán rất tiết kiệm với chỉ \textbf{52{,}48 triệu FLOPs} và \textbf{52{,}822 tham số huấn luyện}, cho phép triển khai hiệu quả trong môi trường thời gian thực mà không cần đến tài nguyên phần cứng lớn.  
Đặc biệt, cơ chế học mặt nạ đặc trưng không chỉ góp phần tối ưu hóa hiệu năng mà còn mang lại khả năng diễn giải nội tại, cho phép xác định rõ những thuộc tính giao dịch nào có ảnh hưởng lớn nhất đến quyết định của mô hình.

Tổng thể, chương 4 đã xác lập một cách thuyết phục rằng MaRENet không chỉ đạt được độ chính xác cao mà còn đảm bảo tính hiệu quả, ổn định và khả năng triển khai thực tế trong các hệ thống giám sát gian lận tài chính hiện đại.










